\section{Finite-block invariants and complete exterior flux}
\label{sec:block}

Fix a real time interval and a stationary simple positively oriented
contour $\Gamma$ whose interior contains exactly $N$ zeros of a
holomorphic heat solution, counted with multiplicity. Assume $H$
does not vanish on a neighborhood of the contour throughout the
interval. Denote the interior divisor by $x_1,\ldots,x_N$ and set
\begin{equation}
 Q(t,z)=\prod_{i=1}^N(z-x_i(t)),\qquad
 R(t,z)=B(t,z)-\frac{Q_z(t,z)}{Q(t,z)}.
 \label{eq:exterior}
\end{equation}
The whole interior divisor has been removed, so $R$ is holomorphic
inside $\Gamma$. It includes the complete exterior and any analytic
factor of $H$; it is not a finite exterior cutoff. Assume also that
the interior contains the block's convex hull, so the center and all
Taylor segments used below lie inside. A cleared enclosing circle
is sufficient.

At a simple interior root the logarithmic derivative has regular
part $\sum_{j\ne i}(x_i-x_j)^{-1}+R(t,x_i)$, so
\begin{equation}
 x_i'=2\sum_{j\ne i}\frac1{x_i-x_j}+2R(t,x_i).
 \label{eq:block-roots}
\end{equation}
Let $c,y_i,p_k$ be the block quantities \eqref{eq:center} and
write $R_i=R(t,x_i)$. Pairwise cancellation, followed by subtraction
of the center velocity, gives
\begin{align}
 c'&=\frac2N\sum_iR_i, \label{eq:center-flux}\\
 p_2'&=2N(N-1)+4\sum_i y_iR_i,\notag\\
 p_3'&=6\sum_i(y_i^2-p_2/N)R_i,\notag\\
 p_4'&=4(2N-3)p_2+8\sum_i(y_i^3-p_3/N)R_i.
 \label{eq:block-powers}
\end{align}
For example, the exterior part of $p_k'$ is
$2k\sum_i(y_i^{k-1}-p_{k-1}/N)R_i$; the internal part is
\eqref{eq:powers}. This proves all lines of
\eqref{eq:block-powers}.

\subsection{Fourth order and the curvature remainder}
For $N\ge4$, retain the polynomial $C_4$ in \eqref{eq:C4} and put
\begin{equation}
 \beta=\frac{2N-3}{N(N-1)},\qquad
 \psi_i=y_i^3-\beta p_2y_i-p_3/N.
 \label{eq:psi}
\end{equation}
Differentiation with respect to a root gives
\begin{equation}
 \partial_{x_i}C_4=-\psi_i,\qquad
 \sum_i\psi_i=0,\qquad \sum_i\psi_iy_i=-4C_4.
 \label{eq:psi-identities}
\end{equation}
The first identity uses
$\partial_{x_i}p_k=k(y_i^{k-1}-p_{k-1}/N)$; the last follows by
substitution of \eqref{eq:C4}. Internal terms cancel in
\eqref{eq:block-powers}, giving the exact balance
\begin{equation}
 C_4'=-2\sum_i\psi_iR_i.
 \label{eq:C4flux}
\end{equation}
Now expand the complete exterior at the current mean:
\begin{equation}
 R(t,c+y)=a(t)+b(t)y+r(t,y),\quad
 a=R(t,c),\quad b=R_z(t,c).
 \label{eq:affine}
\end{equation}
The identities \eqref{eq:psi-identities} yield
\begin{equation}
 C_4'-8bC_4=\E_4,\qquad
 \E_4=-2\sum_i\psi_i r(t,y_i).
 \label{eq:affine-C4}
\end{equation}
In particular an exactly affine exterior gives the constant
$e^{-8\int_{t_0}^t b(v)\dd v}C_4(t)$. For a general exterior,
the same multiplier $\mu(t)$ gives the exact paid balance
\begin{equation}
 \mu(t)C_4(t)-C_4(t_0)=\int_{t_0}^t\mu(v)\E_4(v)\dd v.
 \label{eq:integrated-C4}
\end{equation}

\begin{proposition}[A local curvature allowance]
\label{prop:curvature}
Suppose $|y_i|\le\rho$ and $|R_{zz}|\le C_2$ on every segment
from $c$ to $x_i$. Then
\begin{equation}
 |\E_4|\le C_2\sum_i|\psi_i||y_i|^2\le4NC_2\rho^5.
 \label{eq:curvature}
\end{equation}
\end{proposition}
\begin{proof}
Taylor's integral remainder gives
$r(t,y)=y^2\int_0^1(1-s)R_{zz}(t,c+sy)\dd s$, hence
$|r(t,y)|\le C_2|y|^2/2$. Also
\[
 |\psi_i|\le\rho^3+\beta N\rho^3+\rho^3
          =\left(4-\frac1{N-1}\right)\rho^3\le4\rho^3.
\]
Apply these bounds to \eqref{eq:affine-C4}.
\end{proof}

The fifth power is a local scale allowance. A near-conservation
estimate over a time interval still needs bounds on the initial
value, $b$, $C_2$ and $\rho$ throughout that interval.

\subsection{The general homogeneous law}
\begin{theorem}[Affine-corrected flux for a homogeneous heat integral]
\label{thm:general-flux}
Let $J(x_1,\ldots,x_N)$ be a symmetric polynomial, invariant under
common translation and homogeneous of degree $m\ge1$. Assume it
is an autonomous first integral of the isolated finite heat-root
law \eqref{eq:finite-roots}. Write $J_i=\partial_{x_i}J$ and use
the complete exterior \eqref{eq:exterior}. Then
\begin{align}
 J'&=2\sum_i J_iR_i, \label{eq:Jflux}\\
 J'-2mbJ&=\E_J=2\sum_iJ_i r(t,y_i). \label{eq:Jaffine}
\end{align}
If $Z=z-c$, the single rational kernel
\begin{equation}
 K_J(z)=\sum_i\frac{J_i}{z-x_i}-\frac{mJ}{(z-c)^2}
       =\frac1{Z^2}\sum_i\frac{J_i y_i^2}{Z-y_i}
 \label{eq:Kgeneral}
\end{equation}
satisfies
\begin{equation}
 \E_J=2\av{B(t,z)K_J(z)}.
 \label{eq:contour-general}
\end{equation}
All these symmetric statements extend through internal mergers
while the count and outer clearance persist. For the integrals
$J_j$ of Corollary~\ref{cor:homogeneous}, take $m=j$.
\end{theorem}
\begin{proof}
The chain rule and the assumed isolated first-integral identity
cancel the internal part of \eqref{eq:block-roots}, leaving
\eqref{eq:Jflux}. Translation invariance and Euler's homogeneous
identity give
\begin{equation}
 \sum_iJ_i=0,\qquad \sum_i y_iJ_i=mJ.
 \label{eq:Euler}
\end{equation}
Substitution of \eqref{eq:affine} into \eqref{eq:Jflux} now proves
\eqref{eq:Jaffine}. In particular a purely affine exterior gives
the constant $e^{-2m\int b}J$.

For $w$ inside the contour, Cauchy's formula and the complete
subtraction in \eqref{eq:exterior} give
\begin{equation}
 R(t,w)=\av{\frac{B(t,z)}{z-w}},\qquad
 R_z(t,c)=\av{\frac{B(t,z)}{(z-c)^2}}.
 \label{eq:exterior-Cauchy}
\end{equation}
To verify that the interior poles do not add a term, integrate
$(Q_z/Q)/(z-w)$: all its poles are inside, and it is $O(z^{-2})$
at infinity, so their total residue is zero. The derivative
formula follows by differentiating with respect to $w$. These
statements remain true when $w$ coincides with a block root.
Substitution in $2\sum_iJ_iR_i-2mbJ$ gives
\eqref{eq:contour-general}. Finally,
\[
 \frac1{Z-y_i}=\frac1Z+\frac{y_i}{Z^2}
                   +\frac{y_i^2}{Z^2(Z-y_i)}
\]
and \eqref{eq:Euler} prove the second equality in
\eqref{eq:Kgeneral}. The $Z^{-1}$ and $Z^{-2}$ terms cancel exactly.

For continuation, each raw power sum of the block is
$\av{z^kB}$ and is holomorphic in time on a cleared contour.
Newton identities recover the coefficients of $Q$, and hence every
symmetric polynomial used above. The quotient $H/Q$ is nonvanishing
inside, and its logarithmic derivative $R$ is holomorphic there
and in time. Expressing $J$ as a polynomial in the raw power sums
shows by differentiation that $J_i$ is a polynomial in $x_i$ with
symmetric coefficient functions;
therefore $\sum_iJ_iR(x_i)$ is a finite combination of holomorphic
counted traces. Equivalently, at a multiple slice perturb the heat
solution to $H+\epsilon$, which has simple block zeros for generic
small $\epsilon$, preserves the boundary count and clearance, and
converges with its derivatives there. All displayed symmetric
quantities and first time derivatives converge by their contour
formulas. Passing to the limit proves the identities at the merger,
without assigning a singular individual velocity at that time.
\end{proof}

For $J=C_4$, $J_i=-\psi_i$ and $m=4$. Thus our sign convention
is $K_{C_4}=-\sum_i\psi_i/(z-x_i)-4C_4/(z-c)^2$; it is the
negative of the founding note's $K_4$. The flux
$2\av{BK_{C_4}}$ is exactly the same $\E_4$.
The separate curvature channel is also available from
\begin{equation}
 R_{zz}(t,w)=2\av{\frac{B(t,z)}{(z-w)^3}}.
 \label{eq:Rzz}
\end{equation}
On a circle of radius $A_0$, a uniform source error $\epsilon$
and a distance $\delta$ from $w$ to that circle contribute at most
$2A_0\epsilon/\delta^3$ to this derivative. The next section
retains the cancellations in the combined kernel before taking
absolute values.
